\section{Introduction}
\label{sec:intro}

In the quest to construct general-purpose, multi-task on-device agents, Parameter-Efficient Fine-Tuning (PEFT) techniques, such as Low-Rank Adaptation (LoRA) \cite{hu2021lora}, have emerged as the standard mechanism to adapt massive pre-trained backbones to diverse downstream tasks. By injecting tiny, trainable low-rank adapters alongside frozen base model parameters, researchers can specialize models for tasks ranging from natural language understanding to specialized computer vision with negligible parameter overhead. However, on edge and smart-device deployments, hosting and serving a plethora of task-specific adapters concurrently introduces significant computational, memory, and systems-level scheduling challenges \cite{lane2015deep, dettmers2023qlora}.

To address these serving bottlenecks, the paradigm of \emph{model merging} has received immense interest. Early techniques, such as task arithmetic \cite{ilharco2022editing}, TIES-merging \cite{yadav2023ties}, and DARE \cite{yu2023dare}, attempt to merge several fine-tuned experts directly in the parameter space. While static parameter-space merging is computationally cheap, it forces a single, fixed weight configuration for all incoming queries. Under highly heterogeneous input streams, this static compromise results in representation collapse and severe inter-task interference, particularly when tasks are mutually conflicting.

\begin{figure}[t]
\centering
\includegraphics[width=0.98\columnwidth]{results/fig1.png}
\caption{Joint Mean accuracy of various model merging and routing baselines compared with our proposed \textbf{HyperMerge} under both homogeneous and heterogeneous streams. While Euclidean baselines suffer from representation crowding, HyperMerge achieves a flatline performance of \textbf{83.40\% $\pm$ 5.15\%}, completely immune to stream ordering, operating safely within the logical expert ceiling (99.98\% $\pm$ 0.02\%) and performing on par with state-of-the-art Euclidean baselines with a fully permutation-invariant formulation.}
\label{fig:performance_comparison}
\label{fig1}
\end{figure}

To bypass the limits of static merging, dynamic test-time adapter routing schemes have been proposed. Recent frameworks like Parameter-Free Subspace Routing (PFSR) \cite{pfsr2025} and Sample-wise Activation Blending (SABLE) perform sample-wise, online routing and activation-space blending in a single forward pass. State-of-the-art frameworks, such as Single-Pass Centroid Alignment (SPS-ZCA), utilize zero-shot activation centroids computed during a tiny calibration split to perform dynamic, sample-wise activation ensembling. 

Despite their progress, all existing methods share a common, flat-space geometric foundation: they assume that activation representations and task-specific manifolds lie in a flat Euclidean space ($\mathbb{R}^D$). We argue that this \emph{Euclidean assumption} is fundamentally unsuited for modular deep learning for two primary reasons:
\begin{enumerate}
    \item \textbf{Representation Crowding:} In flat Euclidean space, the volume of a sphere scales polynomially with its dimension ($V(r) \propto r^D$). As a result, when embedding a high number of disjoint task manifolds or hierarchical concepts, representations are forced to crowd around the coordinate origin. This \emph{crowding problem} causes task manifolds to overlap heavily, introducing severe cross-talk (mutual interference) during linear ensembling.
    \item \textbf{Internal Representation and Feature Hierarchy:} Deep neural networks construct representations in a highly hierarchical manner, where early layers encode simple local components (e.g., edges, textures) and deeper layers compose these components into complex semantic taxonomies. This layer-wise progression of abstractions exhibits a natural power-law tree-like scale relationship. Forcing these inherently hierarchical, multi-scale internal feature manifolds into flat Euclidean spaces results in severe topological distortion, which degrades both representation alignment and the precision of distance-based routing.
\end{enumerate}

To shatter these geometric limits, we propose \textbf{HyperMerge} (\textbf{Hyperbolic Space Activation Routing and Fusion}), a radical departure from Euclidean deep learning in model merging. Instead of performing routing and blending in flat planes, HyperMerge maps activation representations and task centroids into the Poincaré Ball model of Hyperbolic Space ($\mathbb{D}_c^D$), which features constant negative curvature. Because the volume of a hyperbolic sphere grows exponentially with its radius ($V(r) \propto e^{(D-1)r}$), hyperbolic space provides a near-infinite ``room'' near its boundary. This negative curvature is mathematically ideal for embedding the power-law tree-like hierarchies of deep representations with minimal distortion, allowing diverse task-expert activations to segregate cleanly and completely eliminating representation crowding and cross-talk.

Using differentiable exponential and logarithmic mapping functions at the origin, HyperMerge projects intermediate activations into the Poincaré Ball in a single forward pass. We introduce \textbf{Hyperbolic Centroid Alignment} (HCA) to compute robust task-specific barycenters (Fréchet means) on a tiny calibration split by projecting Poincaré coordinates to Klein space and computing Einstein midpoints. During online inference, we route queries by measuring hyperbolic geodesic distances to these centroids. We then perform non-linear, permutation-invariant activation ensembling via \textbf{Beltrami-Klein Symmetric Blending} (BKSB), converting Poincaré representations to Klein coordinates to compute symmetric weighted averages (Einstein midpoints) before mapping the merged result back to the Poincaré Ball.

Our empirical evaluation on a 14-layer Analytical Coordinate Sandbox shows that HyperMerge achieves a flatline joint mean accuracy of \textbf{83.40\% $\pm$ 5.15\%} (Figure~\ref{fig:performance_comparison}) under both homogeneous and heterogeneous streams. It completely resolves heterogeneity collapse (0.00\% collapse), performing within the logical bounds of the expert ceiling (99.98\% $\pm$ 0.02\%) and performing on par with the top state-of-the-art Euclidean ensembling baselines but with a mathematically consistent, Lorentz-weighted, and order-independent geometric formulation. 

In summary, our key contributions are:
\begin{itemize}
    \item We challenge the fundamental Euclidean assumption in model merging, demonstrating that flat space introduces representation crowding and destructive cross-talk under heterogeneous deployment.
    \item We present \textbf{HyperMerge}, the first non-Euclidean model merging framework, which leverages Poincaré Ball geometry to natively accommodate internal representation-level hierarchies and cleanly segregate task-expert activations.
    \item We formulate \textbf{Hyperbolic Centroid Alignment} (HCA) based on Einstein midpoints in Klein space, and \textbf{Beltrami-Klein Symmetric Blending} (BKSB) using non-linear Klein algebra to perform permutation-invariant, single-pass, zero-latency dynamic ensembling.
    \item We demonstrate that HyperMerge recovers \textbf{83.40\% $\pm$ 5.15\%} accuracy, completely eliminating representation collapse under heterogeneous streams with perfect, order-independent robustness.
\end{itemize}
